\section{Kesesatan Berpikir Formal (\textit{Logical Fallacies})}

Selain argumen valid, terdapat pola inferensi yang \textbf{tidak sah} namun sering muncul dalam argumen sehari-hari. Dua contoh penting yang menyerupai Modus Ponens dan Modus Tollens---tetapi tidak valid---adalah \textit{fallacy of affirming the consequent} dan \textit{fallacy of denying the antecedent} \cite{rosen2019,iep_natded,libretexts_proofs}.

\subsection{\textit{Fallacy of Affirming the Consequent}}

Kesalahan ini terjadi ketika seseorang menganggap implikasi $p \rightarrow q$ dapat dibalik seolah-olah ekuivalen dengan bikondisional.

\begin{teorema}[Kesalahan: Menetapkan Konsekuen]
Dari $p \rightarrow q$ dan $q$, menyimpulkan $p$ adalah \textbf{tidak valid}.
\end{teorema}

\textbf{Skema yang tidak sah:}
\[
\begin{array}{cl}
p \rightarrow q & \text{(Premis)} \\
q & \text{(Premis)} \\
\hline
\therefore p & \text{(Konklusi tidak sah)}
\end{array}
\]

\begin{contoh}
\begin{itemize}
    \item \textbf{Premis ($p \rightarrow q$)}: ``Jika kompor meledak, maka gedung terbakar.''
    \item \textbf{Fakta ($q$)}: ``Gedung terbakar.''
    \item \textbf{Kesimpulan yang keliru ($\therefore p$)}: ``Pasti kompor yang meledak.''
\end{itemize}
\textbf{Alasan:} Kebakaran dapat disebabkan oleh banyak faktor lain; kebenaran $q$ tidak menjamin kebenaran $p$.
\end{contoh}

\subsection{\textit{Fallacy of Denying the Antecedent}}

Kesalahan ini menyerupai Modus Tollens, tetapi menarik kesimpulan yang salah terhadap konsekuen.

\begin{teorema}[Kesalahan: Menyangkal Anteseden]
Dari $p \rightarrow q$ dan $\neg p$, menyimpulkan $\neg q$ adalah \textbf{tidak valid}.
\end{teorema}

\textbf{Skema yang tidak sah:}
\[
\begin{array}{cl}
p \rightarrow q & \text{(Premis)} \\
\neg p & \text{(Premis)} \\
\hline
\therefore \neg q & \text{(Konklusi tidak sah)}
\end{array}
\]

\begin{contoh}
\begin{itemize}
    \item \textbf{Premis ($p \rightarrow q$)}: ``Jika saya menerima bonus tahun ini, maka tabungan saya bertambah.''
    \item \textbf{Fakta ($\neg p$)}: ``Saya tidak menerima bonus tahun ini.''
    \item \textbf{Kesimpulan yang keliru ($\therefore \neg q$)}: ``Maka tabungan saya tidak dapat bertambah.''
\end{itemize}
\textbf{Alasan:} Tabungan dapat bertambah melalui sumber pendapatan lain; $\neg p$ tidak mengimplikasikan $\neg q$.
\end{contoh}
